\documentclass[11pt]{article}

% parameters are fairly self-explanatory here, thankfully
\usepackage[width=7in,height=9.5in,top=0.75in,centering]{geometry}

% for math-ing
\usepackage{amsmath}
\usepackage{amssymb}


%%%%%% tables %%%%%%%%

\usepackage{array}
\renewcommand{\arraystretch}{1.5}
\setlength{\tabcolsep}{8pt}

%%%%%% plotting %%%%%%

\usepackage[dvipsnames]{xcolor}
\usepackage{tikz}
\usepackage{pgfplots}

\pgfplotsset{every axis/.append style={
                    axis x line=middle,    % put the x axis in the middle
                    axis y line=middle,    % put the y axis in the middle
                    axis line style={<->,color=black}, % arrows on the axis
                    xlabel={$x$},          % default put x on x-axis
                    ylabel={$y$},          % default put y on y-axis
            }}
            
 \pgfplotsset{compat=1.14} % sometimes, everything falls apart without this
 
 %\usepgfplotslibrary{fillbetween} %for shading between curves

%%%%%% commands %%%%%%

% exam heading; course name is first argument, exam information is second argument
\newcommand{\Heading}[2]{
    \fbox{
        {\Large\textbf{#1}}
        }
    \hfill
    \fbox{
        \begin{minipage}{0.2\textwidth}
            \begin{center}
            \textbf{#2}
            \end{center}
        \end{minipage}
    }
    \par\vspace{0.1in}}

% for being lazy while beautifying integrals
\newcommand{\dint}{\displaystyle\int}

% for being lazy while beautifying limits
\newcommand{\dlim}{\displaystyle\lim}

%%%%%% stuff %%%%%%

\usepackage{fancyhdr}

%\setcounter{page}{0} % first page is page 0

\parindent0pt % no indent

\fboxsep=0.25cm % little space in fbox border

%%%%%% BEGIN! %%%%%%%

\begin{document}

\pagestyle{fancy}

% record goals here to create sense of transparency

\fancyhead[c]{%
  \ifcase\value{page}
  % empty test for page = 0
  \or {}% 
  \or
  \or {\bf\color{ProcessBlue} F1: Lines \& Slope}% page=1
  \or {\bf\color{RoyalBlue} L1: Limits and Continuity}% page = 2
  \or {\bf\color{RoyalBlue} L2: Evaluating Limits Using Continuity and Forms}% page = 3
  \or {\bf\color{ForestGreen} D1: Meaning of a Derivative}% page = 4
  \or {\bf\color{ProcessBlue} F2: Exponential and Logarithmic Functions}% page = 5
  \or {\bf\color{ForestGreen} D3: Derivatives of Basic Functions}% page = 7
  \or {\bf\color{ForestGreen} D2: Compute Derivatives}% page = 6
  \or {\bf\color{RedOrange} A1: Derivatives and Graphing}% page = 8
  \or
  \else
  % Empty chead here!
  \fi
}

\thispagestyle{empty} % don't display that this is page 0

\Heading{MATH 1190/1210 Survey of Calculus I}{Knowledge Checkpoint 2\\Standard\\Fall 2025}

\vspace{0.1in}

\textbf{Name} \hrulefill \, \begin{minipage}{0.16\textwidth}\begin{center}\textbf{UVA ID\\ (e.g. hdj4nd)} \end{center}\end{minipage}\hrulefill\\

\begin{minipage}{0.2\textwidth}\begin{center}\textbf{Course \&\\ Section Number} \end{center}\end{minipage} \hrulefill\,\,\textbf{Date} \hrulefill  \\

\hrulefill
\paragraph*{Course \& Section Numbers:}

{\small
\begin{center}
\begin{tabular}{| c | c | c | c | c |}
\hline
%
\begin{minipage}{0.12\textwidth}\begin{center}
1190-100\\
J. Rolf\\
(11:00 AM)
\end{center}\end{minipage}
&
\begin{minipage}{0.12\textwidth}\begin{center}
1190-200\\
J. Rolf\\
(9:30 AM)
\end{center}\end{minipage}
&
\begin{minipage}{0.12\textwidth}\begin{center}
1210-100\\
Mikhail\\
(8 AM)
\end{center}\end{minipage}
&
\begin{minipage}{0.12\textwidth}\begin{center}
1210-200\\
Raul\\
(9:30 AM)
\end{center}\end{minipage}
&
\begin{minipage}{0.12\textwidth}\begin{center}
1210-300\\
Sherman\\
(11 AM)
\end{center}\end{minipage}
\\
%
\hline
%
\begin{minipage}{0.12\textwidth}\begin{center}
1210-400\\
Sherman\\
(12:30 PM)
\end{center}\end{minipage}
&
\begin{minipage}{0.12\textwidth}\begin{center}
1210-500\\
Michael\\
(2 PM)
\end{center}\end{minipage}
&
\begin{minipage}{0.12\textwidth}\begin{center}
1210-600\\
J. Bowling\\
(3:30 PM)
\end{center}\end{minipage}
&
\begin{minipage}{0.12\textwidth}\begin{center}
1210-700\\
Josh\\
(9:30 AM)
\end{center}\end{minipage}
&
\begin{minipage}{0.12\textwidth}\begin{center}
1210-800\\
Mason\\
(2 PM)
\end{center}\end{minipage}
\\
%
\hline
%
\begin{minipage}{0.12\textwidth}\begin{center}
1210-900\\
Caelan\\
(3:30 PM)
\end{center}\end{minipage}
&
\begin{minipage}{0.12\textwidth}\begin{center}
1210-920\\
Caelan\\
(5 PM)
\end{center}\end{minipage}
&
\begin{minipage}{0.12\textwidth}\begin{center}
1210-940\\
Mitchell\\
(10 AM)
\end{center}\end{minipage}
&
\begin{minipage}{0.12\textwidth}\begin{center}
1210-950\\
Mitchell\\
(11 AM)
\end{center}\end{minipage}
&
\begin{minipage}{0.12\textwidth}\begin{center}
1210-960\\
Mitchell\\
(noon)
\end{center}\end{minipage}
\\
\hline
\end{tabular}
\end{center}

}

{%\small

\paragraph*{Instructions:}

\begin{itemize}
    
    \item The regular time allowed for completing this assessment is \fbox{$80$} minutes.
    %\item This assessment is out of 50 points.\\
    \item This assessment contains one single-part or multi-part problem for each of the following targets: {\bf\color{ProcessBlue}F1}, {\bf\color{RoyalBlue}L1}, {\bf\color{RoyalBlue}L2}, {\bf\color{ForestGreen}D1}, {\bf\color{ProcessBlue}F2}, {\bf\color{ForestGreen}D3}, {\bf\color{ForestGreen}D2}, {\bf\color{RedOrange}A1}.
    \item Each problem is worth 10 points.
    %\item You will have opportunities to reassess on the targets appearing on this assessment after the next {\bf 2} assessments.\\
    \item You may NOT use calculators, dictionaries, notes, phones, or any other kind of aid.
    \item Please write neatly. What cannot be read, cannot be graded.
    \item Carefully work out each problem and clearly indicate your final answer to every problem.
    \item Throughout, give the most descriptive answer possible. \textbf{For instance, ``DNE" will not be accepted when ``$=\infty$" is applicable.}
    %\item Please provide the most specific answer possible. For instance, if a limit does not exist because it goes to $\infty$ or $-\infty$, then specify this.
    \item Please use the methods discussed up to this point in the course. For example, any work based on l'H\^opital's Rule will not receive credit.
    \item \textbf{All work must be shown unless otherwise specified.} Partial credit will not be given if appropriate work is not shown.
    \item In particular, {\bf any answer appearing with no justification will receive no credit regardless of whether the answer was correct}, unless otherwise specified.
    \item Please first respond to the prompts on which you feel most proficient, then respond to others later.
    \item This booklet is printed double-sided. It contains 12 pages and 8 numbered problems.
    \end{itemize}
    
    }

\vfill


%%%%%%%%%
\newpage

This page is left blank intentionally.\\

Use this page if you need additional space to write your solutions to any problem. \\

If you wish this page graded, you must refer to this page on \textbf{the original page} of the problem. Otherwise, your grader will NOT consider your work on this page.

\newpage%
%%%%%%%%%


\begin{enumerate}

%%%%%%%%%%%%%%%%%%%%%%%%%%%%%%%%%%%%%%%%%%%%%%%%%%%%%
%%%%%%%%%%% BEGIN PROBLEM F1 %%%%%%%%%%%%%%%%%%%%%%%%
%%%%%%%%%%%%%%%%%%%%%%%%%%%%%%%%%%%%%%%%%%%%%%%%%%%%%

\item The figure below shows a linear function representing a person's distance from their home (in miles) as a function of the amount of time they've spent on a trip (in hours).

\begin{center}
    \includegraphics[scale=0.6]{Fall 2025/KCs/KC2/KC2F1.png}
\end{center}

\begin{enumerate}
    \item How far away was the person from home when they started the trip?

    \vfill

    \item Find the average rate of change of the person's distance from home from the beginning of the trip to 4 hours into the trip based on the graph. To receive full credit, please make sure to include appropriate units.

    \vfill
    \vfill

    \item Use a sentence or two to explain the meaning of the average rate of change found in part (b) \textbf{in the context of this scenario}. 
    
    Note: Simply stating the ``slope" is the ``average rate of change" of the relevant quantities will not receive credit. You must explain the meaning of the slope using the details and quantities related to the context of this question.

    \vfill
    \vfill
\end{enumerate}

%%%%%%%%%%%%%%%%%%%%%%%%%%%%%%%%%%%%%%%%%%%%%%%%%%%%%
%%%%%%%%%%% END PROBLEM F1 %%%%%%%%%%%%%%%%%%%%%%%%%%
%%%%%%%%%%%%%%%%%%%%%%%%%%%%%%%%%%%%%%%%%%%%%%%%%%%%%

\newpage

%%%%%%%%%%%%%%%%%%%%%%%%%%%%%%%%%%%%%%%%%%%%%%%%%%%%%
%%%%%%%%%%% BEGIN PROBLEM L1 %%%%%%%%%%%%%%%%%%%%%%%%
%%%%%%%%%%%%%%%%%%%%%%%%%%%%%%%%%%%%%%%%%%%%%%%%%%%%%

\item For each of the following statements, clearly indicate whether it is true or false. If true, show graphical or verbal example(s) to demonstrate the statement. If false, fully explain why or show a counterexample.

\begin{enumerate}
    \item If $\lim\limits_{x\to -3}f(x)$ DNE, then both $\lim\limits_{x\to -3^-}f(x)$ and $\lim\limits_{x\to -3^+}f(x)$ must also DNE.

    \vfill

    \item If $f(x)$ has a removable discontinuity at $x=a$, then the continuity condition ``$f(a)$ is defined" is never true.

    \vfill
\end{enumerate}

%%%%%%%%%%%%%%%%%%%%%%%%%%%%%%%%%%%%%%%%%%%%%%%%%%%%%
%%%%%%%%%%% END PROBLEM L1 %%%%%%%%%%%%%%%%%%%%%%%%%%
%%%%%%%%%%%%%%%%%%%%%%%%%%%%%%%%%%%%%%%%%%%%%%%%%%%%%

\newpage

%%%%%%%%%%%%%%%%%%%%%%%%%%%%%%%%%%%%%%%%%%%%%%%%%%%%%
%%%%%%%%%%% BEGIN PROBLEM L2 %%%%%%%%%%%%%%%%%%%%%%%%
%%%%%%%%%%%%%%%%%%%%%%%%%%%%%%%%%%%%%%%%%%%%%%%%%%%%%

\item Evaluate the following limits using a theorem about continuity or form.

If you use a theorem about continuity to evaluate a limit, to receive full credit, please provide sufficient explanation for why the theorem can be used. 

If you use form to evaluate a limit, to receive full credit, please list the form and show all steps. 

There's no need to simplify your final answer.

\begin{enumerate}
    \item $\lim\limits_{x\to\infty}\dfrac{x^5+4x^3-3\sqrt{x}}{4\log_2(x)-3x^5}$

    \vfill

    \item $\lim\limits_{x\to 1}\dfrac{x^5+4x^3-3\sqrt{x}}{4\log_2(x)-3x^5}$

    \vfill
\end{enumerate}

%%%%%%%%%%%%%%%%%%%%%%%%%%%%%%%%%%%%%%%%%%%%%%%%%%%%%
%%%%%%%%%%% END PROBLEM L2 %%%%%%%%%%%%%%%%%%%%%%%%%%
%%%%%%%%%%%%%%%%%%%%%%%%%%%%%%%%%%%%%%%%%%%%%%%%%%%%%

\newpage

%%%%%%%%%%%%%%%%%%%%%%%%%%%%%%%%%%%%%%%%%%%%%%%%%%%%%
%%%%%%%%%%% BEGIN PROBLEM D1 %%%%%%%%%%%%%%%%%%%%%%%%
%%%%%%%%%%%%%%%%%%%%%%%%%%%%%%%%%%%%%%%%%%%%%%%%%%%%%

\item Sam drives from home to school daily, where the commute takes approximately 20 minutes.

Below is a graph representing the relationship between the distance $s$ Sam has driven (in miles) and time $t$ (in minutes) during the commute.

In case it's unclear, there's no corner in the graph.

\begin{center}
        \includegraphics[scale=0.5]{Fall 2025/KCs/KC2/KC2D1.png}
    \end{center}

\begin{enumerate}
    \item Which one of i, ii, and iii best describes Sam's drive? Explain why.
    
    To receive full credit, point out at least two \textbf{graphical features related to derivative} that matches the description.

    \begin{enumerate}
        \item Sam drove on city roads where they stopped for a red light approximately every 5 minutes.
        \item Sam started driving, stopped at a stop sign in the neighborhood, then continued on a main road. They stopped again at a red light right before getting onto the highway. They drove at a constant speed on the highway for a 5-10 minutes before slowing down and eventually stopped as they arrived at the school.
        \item Sam started driving on a highway with a high speed for about 10 minutes, but then had to reduce the speed due to construction. Once they left the construction zone, they smoothly accelerated to a faster speed and drove at that speed for before arriving at the school and gradually come to a full stop.
    \end{enumerate}

    \vfill

    \item Rank the values of $s'(10)$, $s'(14)$, and $s'(19)$ from smallest to largest. No explanation required.

    \vspace{0.5cm}
    
    \begin{center}
    \large{\rule[-8pt]{2cm}{0.4pt}$<$\rule[-8pt]{2cm}{0.4pt}$<$\rule[-8pt]{2cm}{0.4pt}}
    \end{center}
\end{enumerate}

%%%%%%%%%%%%%%%%%%%%%%%%%%%%%%%%%%%%%%%%%%%%%%%%%%%%%
%%%%%%%%%%% END PROBLEM D1 %%%%%%%%%%%%%%%%%%%%%%%%%%
%%%%%%%%%%%%%%%%%%%%%%%%%%%%%%%%%%%%%%%%%%%%%%%%%%%%%

\newpage

%%%%%%%%%%%%%%%%%%%%%%%%%%%%%%%%%%%%%%%%%%%%%%%%%%%%%
%%%%%%%%%%% BEGIN PROBLEM F2 %%%%%%%%%%%%%%%%%%%%%%%%
%%%%%%%%%%%%%%%%%%%%%%%%%%%%%%%%%%%%%%%%%%%%%%%%%%%%%

\item Respond to the following prompts. For this specific question, you are not required to compute any powers of constants or logarithms. For example, if your final answer includes $10^3$ or $\log_2(5)$, you can leave it as is.

\begin{enumerate}
    \item Simplify the expression below as much as possible so that the final result contains only positive exponents. You may expand exponential expressions and/or apply properties of exponents as needed. 

    \[\left(\dfrac{3^{-5/2}x}{\sqrt{3}x^{-3/5}}\right)^5\]


    \vfill

    \item Solve the following equation for $x$. 

    \[7^{2x-5}=10\]

    \vfill
\end{enumerate}


%%%%%%%%%%%%%%%%%%%%%%%%%%%%%%%%%%%%%%%%%%%%%%%%%%%%%
%%%%%%%%%%% END PROBLEM F2 %%%%%%%%%%%%%%%%%%%%%%%%%%
%%%%%%%%%%%%%%%%%%%%%%%%%%%%%%%%%%%%%%%%%%%%%%%%%%%%%

\newpage

%%%%%%%%%%%%%%%%%%%%%%%%%%%%%%%%%%%%%%%%%%%%%%%%%%%%%
%%%%%%%%%%% BEGIN PROBLEM D3 %%%%%%%%%%%%%%%%%%%%%%%%
%%%%%%%%%%%%%%%%%%%%%%%%%%%%%%%%%%%%%%%%%%%%%%%%%%%%%

\item Write down the derivatives of the following functions. No work required. No need to simplify.

If you need to write down your work, please do so at the bottom half of this page, not in the answer space.

\begin{center}
        \begin{tabular}{l c r l c r}
            %
            $\displaystyle\frac{d}{dx}\left[ 9x^{6} \right]$ & $=$ & \rule[-8pt]{1.5in}{0.4pt} &
            %
            $\displaystyle\frac{d}{dx}\left[ \dfrac{1}{\sqrt[5]{x^3}} \right]$ & $=$ & \rule[-8pt]{1.5in}{0.4pt}\\[0.3in]
            %%%%%
            $\displaystyle\frac{d}{dx}\left[ \left(\dfrac{3}{4}\right)^x \right]$ & $=$ & \rule[-8pt]{1.5in}{0.4pt} & 
            %%%%%
            $\displaystyle\frac{d}{dx}\left[ \pi^{e} \right]$ & $=$ & \rule[-8pt]{1.5in}{0.4pt}\\[0.3in]
            %%%%%
            $\displaystyle\frac{d}{dx}\left[ x^5+3 \right]$ & $=$ & \rule[-8pt]{1.5in}{0.4pt}&
            %
            $\displaystyle\frac{d}{dx}\left[ \sqrt{x}-\ln(x) \right]$ & $=$ & \rule[-8pt]{1.5in}{0.4pt}\\[0.3in]
            %%%%%
            $\displaystyle\frac{d}{dx}\left[ \dfrac{x^3+5x^2-x}{x^4} \right]$ & $=$ & \rule[-8pt]{1.5in}{0.4pt} & 
            %
            $\displaystyle\frac{d}{dx}\left[ \log_{2}(x) \right]$ & $=$ & \rule[-8pt]{1.5in}{0.4pt}\\[0.3in]
            %%%%%
        \end{tabular}
        
\end{center}


%%%%%%%%%%%%%%%%%%%%%%%%%%%%%%%%%%%%%%%%%%%%%%%%%%%%%
%%%%%%%%%%% END PROBLEM D3 %%%%%%%%%%%%%%%%%%%%%%%%%%
%%%%%%%%%%%%%%%%%%%%%%%%%%%%%%%%%%%%%%%%%%%%%%%%%%%%%

\newpage

%%%%%%%%%%%%%%%%%%%%%%%%%%%%%%%%%%%%%%%%%%%%%%%%%%%%%
%%%%%%%%%%% BEGIN PROBLEM D2 %%%%%%%%%%%%%%%%%%%%%%%%
%%%%%%%%%%%%%%%%%%%%%%%%%%%%%%%%%%%%%%%%%%%%%%%%%%%%%

\item Find the derivative of the following expressions. 

A correct final answer receives full credit. Nevertheless, you are encouraged to write down the steps so that you will be able to receive partial credit if you final answer is not fully correct. There's no need to simplify your final answer.

\begin{enumerate}
\item $\dfrac{\ln(x)+x^2}{x^3-2x}$

\vfill

\item $4\sqrt{x}\left(e^2+3x^4\right)^2$

\vfill 
\end{enumerate}

%split part (a) and part (b)

%%%%%%%%%%%%%%%%%%%%%%%%%%%%%%%%%%%%%%%%%%%%%%%%%%%%%
%%%%%%%%%%% END PROBLEM D2 %%%%%%%%%%%%%%%%%%%%%%%%%%
%%%%%%%%%%%%%%%%%%%%%%%%%%%%%%%%%%%%%%%%%%%%%%%%%%%%%

\newpage

%%%%%%%%%%%%%%%%%%%%%%%%%%%%%%%%%%%%%%%%%%%%%%%%%%%%%
%%%%%%%%%%% BEGIN PROBLEM A1 %%%%%%%%%%%%%%%%%%%%%%%%
%%%%%%%%%%%%%%%%%%%%%%%%%%%%%%%%%%%%%%%%%%%%%%%%%%%%%

\item Consider the following function $f(x)$.

\[f(x)=x^3-6x^2+5\]

Respond to the following prompts. Make sure to show all work.

\begin{enumerate}
    \item Find all critical numbers of $f(x)$.

    \vfill

    \item On what interval(s) is $f(x)$ concave up? On what intervals(s) is $f(x)$ concave down?

    \vfill
\end{enumerate}


%%%%%%%%%%%%%%%%%%%%%%%%%%%%%%%%%%%%%%%%%%%%%%%%%%%%%
%%%%%%%%%%% END PROBLEM A1 %%%%%%%%%%%%%%%%%%%%%%%%%%
%%%%%%%%%%%%%%%%%%%%%%%%%%%%%%%%%%%%%%%%%%%%%%%%%%%%%

\newpage

This page is left blank intentionally.\\

Use this page if you need additional space to write your solutions to any problem. \\

If you wish this page graded, you must refer to this page on \textbf{the original page} of the problem. Otherwise, your grader will NOT consider your work on this page.

\newpage

This page is left blank intentionally.\\

Use this page if you need additional space to write your solutions to any problem. \\

If you wish this page graded, you must refer to this page on \textbf{the original page} of the problem. Otherwise, your grader will NOT consider your work on this page.

\end{enumerate}

\end{document}